% =============================================================================
% Cap. 3 — Implementação (estado actual à data do intercalar, conforme guia §3)
% =============================================================================

\section{Implementação --- Estado actual}
\label{cap:implementacao}

Este capítulo documenta o estado de implementação à data deste
relatório intercalar (3~de~Maio~de~2026), apresentando a
decomposição de tarefas com precedências (\S\ref{sec:tasks}), o
estado por componente (\S\ref{sec:components}), o cumprimento dos
critérios de aceitação (\S\ref{sec:acceptance}), os desvios ao plano
(\S\ref{sec:deviations}) --- incluindo o reconhecimento honesto da
demo interna síncrona não realizada na janela prevista ---,
a demonstração em produção (\S\ref{sec:demo}), o estado dos testes
automatizados (\S\ref{sec:tests-state}) e o calendário actualizado
para a fase restante (\S\ref{sec:calendario-restante}).

\subsection{Decomposição do trabalho}
\label{sec:tasks}

O trabalho foi decomposto em onze tarefas principais com precedências
implícitas, conforme a Tabela~\ref{tab:tasks}.

\begin{longtable}{@{}cL{6.0cm}L{4.5cm}L{2.0cm}@{}}
\toprule
\textbf{T} & \textbf{Tarefa} & \textbf{Precedências} & \textbf{Estado} \\
\midrule
\endhead
T1  & Configuração do repositório, proposta e MoSCoW            & ---       & Concluída \\
T2  & ADRs 0001--0010 com alternativas consideradas             & T1        & Concluída \\
T3  & Modelo de dados core                                      & T2        & Concluída \\
T4  & Autenticação JWT em \emph{cookies} HttpOnly + CSRF        & T3        & Concluída \\
T5  & API REST e \emph{viewsets}                                & T3, T4    & Concluída \\
T6  & \emph{Frontend} \emph{vanilla} \emph{mobile-first}        & T5        & Concluída \\
T7  & Imutabilidade tripla (ORM, DRF, \emph{triggers} PG)       & T3, T5    & Concluída \\
T8  & Auditoria interna de segurança e correcções               & T7        & Concluída \\
T9  & \emph{Deployment} Fly.io + Neon; HSTS \emph{preload}      & T6        & Concluída \\
T10 & \emph{Wizard} de evidência, \emph{lookups}, exportação PDF & T5, T6   & Concluída \\
T11 & Relatório intercalar                                      & T1--T10   & Em curso \\
\bottomrule
\caption{Decomposição de tarefas com precedências.}
\label{tab:tasks}
\end{longtable}

\subsection{Estado final por componente}
\label{sec:components}

A Tabela~\ref{tab:components} resume o estado de cada componente. O
estado verde indica conclusão validada por testes; amarelo indica
funcionalidade entregue com limitação reconhecida; vermelho indica
funcionalidade fora do âmbito do MVP ou ainda em curso.

\begin{longtable}{@{}L{4cm}L{2cm}L{8cm}@{}}
\toprule
\textbf{Componente} & \textbf{Estado} & \textbf{Observação} \\
\midrule
\endhead
\emph{Backend} Django + DRF                  & \cellcolor{green!12}\textcolor{successgreen}{verde}  & 213 testes a passar; modelos, \emph{viewsets}, serializadores, permissões. \\
\midrule
Cadeia de custódia (SHA-256 encadeado)       & \cellcolor{green!12}\textcolor{successgreen}{verde}  & \texttt{transaction.atomic} + \texttt{select\_for\_update} (correcções da auditoria interna). \\
\midrule
Imutabilidade tripla                         & \cellcolor{green!12}\textcolor{successgreen}{verde}  & ORM + DRF + \emph{triggers} PostgreSQL (\emph{migration} 0002). \\
\midrule
Autenticação JWT em \emph{cookies}           & \cellcolor{green!12}\textcolor{successgreen}{verde}  & HttpOnly + Secure + SameSite=Strict; CSRF \emph{double-submit}; \emph{rate limit} 5/min. \\
\midrule
\emph{Frontend} \emph{vanilla}               & \cellcolor{green!12}\textcolor{successgreen}{verde}  & 13 vistas \emph{mobile-first}; sem \emph{build step}; tema dia/noite com sunset NOAA. \\
\midrule
Mapa Leaflet + OpenStreetMap                 & \cellcolor{green!12}\textcolor{successgreen}{verde}  & Leaflet 1.9.4 com SRI SHA-512; Referrer-Policy \texttt{strict-origin-when-cross-origin}. \\
\midrule
Exportação PDF (RF6)                         & \cellcolor{green!12}\textcolor{successgreen}{verde}  & ReportLab; \texttt{/api/evidences/<id>/pdf/} e \texttt{/api/occurrences/<id>/pdf/} com declaração ISO/IEC~27037. \\
\midrule
Exportação CSV                               & \cellcolor{green!12}\textcolor{successgreen}{verde}  & \emph{Streaming} com \emph{cap} de 10 000 linhas; respeita \emph{ownership}; \emph{audit log} regista cada exportação. \\
\midrule
Lookup IMEI/VIN                              & \cellcolor{green!12}\textcolor{successgreen}{verde}  & \texttt{DatabaseCache} 30 dias (IMEI) e 180 dias (VIN). \\
\midrule
\emph{Deployment}                            & \cellcolor{green!12}\textcolor{successgreen}{verde}  & Fly.io fra + Neon; HTTPS A+ (Qualys); HSTS \emph{preload} submetido; Mozilla Observatory A+. \\
\midrule
Cobertura de testes                          & \cellcolor{orange!15}\textcolor{warningorange}{amarelo} & 67{,}4\% global; áreas com cobertura insuficiente reconhecidas (\texttt{validators.py}, \texttt{imei\_lookup.py}, \texttt{seed\_demo.py}). \\
\midrule
\emph{Backups} de \texttt{media/}            & \cellcolor{red!15}\textcolor{dangerred}{vermelho}    & Volume Fly.io persistente sem \emph{backup} externo automatizado; trabalho futuro. \\
\midrule
SAST/SCA/DAST em CI                          & \cellcolor{red!15}\textcolor{dangerred}{vermelho}    & Fora do âmbito do MVP; documentado para Fase~3. \\
\bottomrule
\caption{Estado por componente à data do intercalar.}
\label{tab:components}
\end{longtable}

\subsection{Cumprimento dos critérios de aceitação}
\label{sec:acceptance}

A Tabela~\ref{tab:acceptance} confronta o estado dos requisitos
funcionais (RF) e não-funcionais (RNF) inscritos em
\texttt{docs/scope/requirements.tex} (v1.0) com a sua verificação
empírica à data deste intercalar. Os requisitos \emph{Must have}
foram todos cumpridos. Os \emph{Should have} foram cumpridos. Os
\emph{Could have} estão parcialmente cumpridos. Os \emph{Won't have}
permaneceram fora do âmbito.

\begin{longtable}{@{}cL{8.5cm}L{2.0cm}L{2.5cm}@{}}
\toprule
\textbf{ID} & \textbf{Requisito} & \textbf{Prioridade} & \textbf{Estado} \\
\midrule
\endhead
RF01 & Autenticação JWT com perfis                       & Must  & \cellcolor{green!12}Cumprido \\
RF02 & Ocorrências com georreferenciação                 & Must  & \cellcolor{green!12}Cumprido \\
RF03 & Registo de prova com foto, GPS, \emph{hash}       & Must  & \cellcolor{green!12}Cumprido \\
RF04 & Ficha de dispositivo digital                       & Must  & \cellcolor{green!12}Cumprido \\
RF05 & Cadeia de custódia \emph{append-only} encadeada    & Must  & \cellcolor{green!12}Cumprido \\
RF06 & Exportação PDF                                     & Must  & \cellcolor{green!12}Cumprido \\
RF07 & API REST com 10+ \emph{endpoints} documentados     & Must  & \cellcolor{green!12}Cumprido (mais de 12 \emph{endpoints}) \\
RF08 & \emph{Dashboard} por estado                        & Should & \cellcolor{green!12}Cumprido \\
RF09 & Pesquisa e filtragem                               & Should & \cellcolor{green!12}Cumprido \\
RF10 & Perfis adicionais (coordenador, magistrado)        & Could & \cellcolor{orange!15}Estrutura preparada (ADR-0006) \\
RF11 & Verificação de IMEI via \emph{API} externa         & Could & \cellcolor{green!12}Cumprido (ADR-0008) \\
RF12 & \emph{Dashboard} analítico                         & Could & \cellcolor{green!12}Cumprido (\texttt{/stats/}) \\
RF13 & Despachos judiciais                                & Could & \cellcolor{red!15}Não realizado \\
RF14--RF17 & Integração PSP, dados reais, módulos QBF, \emph{app} nativa & Won't & \cellcolor{gray!15}Fora do âmbito \\
\midrule
RNF01 & HTTPS obrigatório                                 & Must  & \cellcolor{green!12}Cumprido (HSTS \emph{preload} + Qualys A+) \\
RNF02 & \emph{Mobile-first} utilizável com uma só mão      & Must  & \cellcolor{green!12}Cumprido (WCAG 2.1 AA) \\
RNF03 & Integridade SHA-256 e \emph{append-only}           & Must  & \cellcolor{green!12}Cumprido \\
RNF04 & Resposta inferior a 2 s                           & Should & \cellcolor{green!12}Verificado em produção \\
RNF05 & Testes automatizados em módulos críticos           & Should & \cellcolor{green!12}213 testes a passar \\
RNF06 & \emph{Deployment} público                         & Could & \cellcolor{green!12}\href{https://forensiq.pt}{forensiq.pt} \\
\bottomrule
\caption{Cumprimento dos requisitos funcionais e não-funcionais.}
\label{tab:acceptance}
\end{longtable}

\subsection{Desvios ao plano}
\label{sec:deviations}

A execução afastou-se em três pontos do plano apresentado na
proposta; todos foram comunicados ao orientador.

\paragraph{\emph{Stack} do \emph{backend}.}
A proposta indicava \emph{``Django + DRF''} como opção principal,
sem alternativa. A escolha foi confirmada em ADR-0002 e mantida.
Sem desvio.

\paragraph{Demo interna síncrona não realizada na janela prevista
(Sem.~7).}
O calendário do orientador previa demo interna síncrona via Teams
entre 28~de~Abril e 2~de~Maio de 2026 (Sem.~7). A demo síncrona
\emph{não foi realizada} nessa janela. O autor reconhece o desvio.
Como mitigação, propõe-se ao orientador, em comunicação paralela a
este relatório, que a demonstração seja feita de forma \emph{assíncrona}
através do site em produção \href{https://forensiq.pt}{forensiq.pt},
com credenciais de demonstração entregues por canal separado. Esta
modalidade permite ao orientador exercitar todos os fluxos críticos
(criação de ocorrência, registo de evidência com GPS e fotografia,
transição de custódia, exportação de PDF) sem necessidade de
configurar ambiente local. Caso o orientador prefira demo síncrona,
o autor está disponível para a agendar nas semanas 9~ou~10
(7--16~mai), com a vantagem de já ter incorporado o feedback do
intercalar.

\paragraph{Auditorias internas extras (16, 18 e 19~Abr).}
O plano original avançava da implementação directamente para o
relatório intercalar. Em meados de Abril realizaram-se três
auditorias internas (segurança, \emph{design}, taxonomia) que
detectaram pontos a corrigir, em particular uma \emph{race
condition} na cadeia de custódia (\emph{fix~B-C3}) e a necessidade
de tornar o cálculo do \emph{hash} puro (\emph{fix~B-C2}). O esforço
adicional consumiu cerca de uma semana mas foi essencial para a
robustez actual.

\subsection{Demonstração em produção}
\label{sec:demo}

O ambiente de produção em \href{https://forensiq.pt}{forensiq.pt} é
uma cópia plenamente funcional do ForensiQ, alojada em Fly.io
Frankfurt com PostgreSQL gerido pela Neon na mesma região
(\S\ref{sec:c4-l2}). Está disponível um conjunto de demonstração
realista, gerado pelo \emph{management command}
\texttt{seed\_demo} (\texttt{src/backend/core/management/commands/seed\_demo.py}),
com cinco ocorrências em cidades portuguesas (Lisboa, Porto,
Coimbra, Braga, Faro), respectivos NUIPCs e treze itens de prova
distribuídos pelos vários tipos da taxonomia, incluindo cartões SIM
como sub-componentes hierárquicos. As fotografias são imagens
\emph{placeholder} sintéticas geradas pelo \emph{seed} para evitar
qualquer dado pessoal real.

O acesso à demonstração é feito com três contas geradas no momento
do \emph{seed} com \emph{passwords} aleatórias (16 caracteres):
\texttt{agente.demo} (perfil \texttt{AGENT}, badge GNR-12345),
\texttt{perito.demo} (perfil \texttt{EXPERT}, badge PJ-LPC-007) e
\texttt{pedro.pestana} (perfil \texttt{EXPERT} com \emph{flag}
\texttt{is\_staff}, badge UAB-ORIENT). As \emph{passwords} são
entregues ao orientador em comunicação paralela ao envio deste
relatório, e devem ser rodadas por ele ao primeiro acesso.

\subsection{Estado dos testes automatizados}
\label{sec:tests-state}

À data deste intercalar, a \emph{suite} automatizada do ForensiQ
contém \textbf{213 testes a passar} sob \texttt{pytest 9.0.3}
(execução em \(\texttt{2026{-}05{-}03}\), tempo total
\(\texttt{0:02:47}\)). A distribuição por \emph{suite} é:
\texttt{tests.py} (14 testes de modelos), \texttt{tests\_api.py}
(83 testes de API: autenticação, IDOR, imutabilidade, transições,
\emph{throttle}, \emph{lookups}, estatísticas, CSP, \emph{rate
limiting}), \texttt{tests\_pdf.py} (18 testes de exportação PDF),
\texttt{tests\_frontend.py} (54 testes de vistas e \emph{templates}),
\texttt{tests\_new\_features.py} (22 testes de funcionalidades
introduzidas em iterações posteriores) e \texttt{tests\_table\_mode.py}
(22 testes de modo tabela e exportação CSV).

A cobertura medida por \texttt{coverage.py} é de \textbf{67{,}4\%}
global em \texttt{core/} (1881 \emph{statements}, 536 não cobertos).
A leitura honesta destes números é a seguinte: a cobertura
concentra-se nas áreas críticas para integridade forense ---
\texttt{models.py} (78{,}9\%), \texttt{views.py} (75{,}1\%),
\texttt{pdf\_export.py} (86{,}7\%), \texttt{middleware.py} (90{,}7\%) ---;
a baixa cobertura global é puxada por três módulos:
\texttt{seed\_demo.py} (\emph{script} de demonstração executado
manualmente, intencionalmente sem testes), \texttt{imei\_lookup.py}
(integração externa que exigiria \emph{mocking} mais sofisticado) e
\texttt{validators.py} (cujos casos felizes estão cobertos
indirectamente pelos testes de criação de evidência). O reforço da
cobertura é uma das prioridades das semanas 9--10
(\S\ref{sec:calendario-restante}).

\subsection{Calendário actualizado para a fase restante}
\label{sec:calendario-restante}

A Tabela~\ref{tab:calendario-restante} apresenta o calendário
actualizado para a fase restante, da semana seguinte ao intercalar
até à entrega do relatório final em 24~de~Junho de 2026.

\begin{longtable}{@{}L{1cm}L{2.4cm}L{11.0cm}@{}}
\toprule
\textbf{Sem.} & \textbf{Datas} & \textbf{Trabalho previsto} \\
\midrule
\endhead
9--10 & 7--16 mai & Reforço de cobertura de testes (alvo $\geq 75\%$); \emph{property-based testing} de validadores; eventual demo síncrona com o orientador. \\
\midrule
11--12 & 19--30 mai & Polimento da \emph{interface} \emph{mobile-first}; tabela de auditoria visual no \emph{dashboard}; preparação dos anexos do relatório final (manifesto de testes, capturas de ecrã). \\
\midrule
13 & 2--6 jun & Revisão geral; validação cruzada dos critérios de aceitação; congelamento do âmbito. \\
\midrule
14 & 9--13 jun & Redacção dos Capítulos~4 (Testes) e~5 (Conclusões) do relatório final; bibliografia APA consolidada. \\
\midrule
15 & 16--20 jun & Reunião de preparação para defesa com o orientador (síncrono); ensaio de perguntas de júri. \\
\midrule
16 & 24 jun & Submissão do relatório final e do código. \\
\bottomrule
\caption{Calendário actualizado para a fase restante.}
\label{tab:calendario-restante}
\end{longtable}
